\section{Discussion}\label{sec:disc}

\subsection{What the Design Buys and What It Costs}
Compared with tokens and CRLs, the central gain is the removal of an implicit-trust window: on the audited path a revoked device loses access at the next call, and on the fast path within the replica lag, a bound the operator can set and measure (Propositions~\ref{prop:revoke} and~\ref{prop:stale}; 33--147 accepted requests for the baselines against none or $\rho\Delta$). Quorum governance removes the single-registrar insider (T7) at a gas cost that scales with the voter group rather than the fleet. The tiered policy concentrates verification effort where a wrong decision hurts: at T2--T3 the gate and strict mode cut a spoofing attacker's acceptance from about 100\% to 0\% and from about half to zero, and the cost is a few percent of legitimate requests denied when the process is upset. In contrast to schemes that authenticate once and trust the session \cite{wang2026medicaliot}, nothing here is trusted by default; in contrast to conceptual smart-factory designs \cite{farhan2026ztfsf}, the policy is an executable contract with measured cost.

The costs are equally concrete. The fast path doubles per-request compute (two signature checks), the audited path is bounded by block gas (about 426 decisions per 30M-gas block by gas, and only 11--18 per second measured on our single-host four-validator network), recall of very large fleets takes tens of blocks, and the whole design adds a consortium chain, a replica and an oracle to operate. For a plant that has one owner and a good internal PKI, this is probably unjustified; the argument for it is multi-party plants in which revocation authority and audit trails must be shared.

\subsection{Deployment Guidance}
\emph{Tier assignment.} Tiers should follow consequence, not protocol: any action whose misuse could damage equipment or product, or alter safety-relevant logic, belongs on the audited, strict path; high-rate reads belong on the fast path. \emph{Gate calibration.} The anomaly gate should be set from a plant's own distribution of legitimate deviations, including start-up and batch changes; Section~\ref{sec:gate_eval} shows the operating point moves substantially with the upset rate. \emph{Replica lag.} $\Delta$ should be measured, monitored and alarmed, since Proposition~\ref{prop:stale} is only as good as A5. \emph{Recall rehearsal.} Batch revocation should be rehearsed with quarantine (\texttt{suspendBatch}) first, because REVOKED is terminal and a mistaken revocation requires re-enrolment with a new key. \emph{Gas economics.} On a private PoA network gas can be zero-priced, so the figures above indicate capacity and relative cost rather than money; gasless relaying can hide fees from constrained devices \cite{madhwal2025gasless}.

\subsection{Threats to Validity}
\emph{Internal validity.} In the main experiments, devices, signal ranges and attacks are synthetic and chosen by us, and the legitimate and attacking anomaly ranges overlap by construction. In the TON\_IoT study only the anomaly signal is real: posture, network context and reliability history are synthetic, the assignment of a flow score to a device is arbitrary, attackers are assumed to spoof posture and network context perfectly, and the held-out-type test approximates but is not a real zero-day setting. The dataset itself has known weaknesses that we saw directly: exact duplicates, collection of normal and attack records on different days, an encoding artefact in one file, and attacks from a few source hosts, so the supervised AUC of 0.9999 is optimistic. None of this is a plant measurement. The baselines are minimal logical-clock models: short-lived certificates and OCSP stapling are modelled only by their validity window, and a production PKI would add renewal failures, clock skew and caching. The outage result is policy-defined. \emph{External validity.} Everything ran on one two-vCPU container with the two tracks using different signature schemes; no measurement was made on a microcontroller or on industrial protocols, and the four-validator Besu network ran on that same host over the loopback interface, so it has no wide-area latency, its throughput is limited by CPU contention, and the fault tests used a single pending transaction and small samples. \emph{Construct validity.} In the simulation, revocation exposure is counted on a logical clock and ignores propagation and finality time; the Besu measurements give the finality component (one block period) but not wide-area propagation. \emph{Conclusion validity.} Intervals come from 10--30 runs on one host; the 500-trial attack tests have a 0.6\% rule-of-three bound.

\subsection{Limitations and Future Work}
First, the multi-node measurements come from a single host; a deployment across several machines and sites is needed to measure wide-area replication lag $\Delta$, realistic throughput and behaviour under partitions. Second, the PEP is trusted for request-signature checking in Track A, and a single oracle and registrar remain single points of compromise (A4); a threshold-signed oracle \cite{hassan2023oracle} and a quorum registrar set would remove them. Third, our legacy-device proxy is a design with a modelled class, not an implemented gateway, and no real fieldbus or OPC-style protocol was exercised. Fourth, trust scores and anomaly values appear in transaction data and are visible to consortium members; differential-privacy protection of exposed scores \cite{ren2026ditdp} is not implemented. Fifth, adversarial-AI mimicry is richer than our populations \cite{usama2026multilayer}, and the signal producers (attestation, process-aware anomaly detection, possibly cross-site learning \cite{tom2025flcti}) are inputs we assume. Sixth, multi-PEP zones require binding the PEP identity into requests (Proposition~\ref{prop:replay}). Seventh, the policy parameters were hand-tuned, and learning them from plant data \cite{salim2026ztblociot} is open. Eighth, the real-data study shows that a detector-dependent gate inherits the detector's blind spots on unseen attacks and is exposed to an adaptive attacker who can steer its own score; adversarially robust detection, rate limiting of low-score probing and joint use of several detectors are the natural defences, and we did not evaluate them. Ninth, a field study is missing: per-device telemetry joined to the traffic of the same device, with real attestation, would replace the arbitrary pairing used here. Finally, the contracts are research code: they pass 16 tests, including a randomised state-machine test, and a Slither run found no serious issue after one hardening fix, but they have not been audited or formally verified, and the literature on smart-contract quality assurance \cite{wei2024quality} should be followed before any production use.
